\documentclass[12pt,a4paper]{article}

\usepackage[a4paper,margin=1in]{geometry}
\usepackage{graphicx}
\usepackage{booktabs}
\usepackage{amsmath,amssymb}
\usepackage{float}
\usepackage{hyperref}
\usepackage{xcolor}
\usepackage{listings}
\usepackage{enumitem}
\usepackage{fancyhdr}
\usepackage{longtable}

\hypersetup{colorlinks=true,linkcolor=blue,urlcolor=blue}

\pagestyle{fancy}
\fancyhf{}
\lhead{ICS1512}
\rhead{Machine Learning Algorithms Laboratory}
\cfoot{\thepage}

\lstset{
language=Python,
basicstyle=\ttfamily\small,
numbers=left,
frame=single,
breaklines=true,
showstringspaces=false,
keywordstyle=\color{blue},
commentstyle=\color{gray},
stringstyle=\color{purple}
}

\begin{document}

\begin{tabular}{|p{4cm}|p{11cm}|}
\hline
Experiment No. & 4 \\ \hline
Experiment Title & Binary Classification using Linear and Kernel-Based Models
\footnote{\textbf{GitHub Repository: }\url{https://github.com/Nightingales21/ICS1512-Machine-Learning-Labratory}}\\ \hline
Student Name & Nitin \\ \hline
Register Number & 312224700142 \\ \hline
Faculty & Dr. Poreddy Ajay Kumar Reddy \\ \hline
Submission Date & \today \\ \hline
\end{tabular}

\section{Objective}
To classify emails as spam or ham using Logistic Regression and Support
Vector Machine (SVM) classifiers, and to analyze the effect of
hyperparameter tuning -- regularization strength and type for Logistic
Regression, and kernel choice, $C$, and $\gamma$ for SVM -- on
classification performance.

\section{Dataset Description}
\begin{longtable}{|p{6cm}|p{8cm}|}
\hline
Dataset Name & Spambase \\ \hline
Dataset Source & UCI Machine Learning Repository / Kaggle
(\url{https://www.kaggle.com/datasets/somesh24/spambase}) \\ \hline
Number of Samples & 4{,}601 \\ \hline
Number of Features & 57 (48 word-frequency, 6 character-frequency, 3
capital-run-length features) \\ \hline
Number of Classes & 2 (Spam = 1, Ham = 0) \\ \hline
Missing Values & 0 (verified programmatically) \\ \hline
Train--Test Split & 80\% / 20\%, stratified, \texttt{random\_state=42} \\ \hline
\end{longtable}

Class balance: 2{,}788 ham emails (60.6\%) and 1{,}813 spam emails (39.4\%).
This is the same dataset used in Experiment 2, allowing a direct
cross-experiment comparison between the generative/instance-based models
studied there (Naive Bayes, KNN) and the linear/margin-based models studied
here (Logistic Regression, SVM).

\section{Preprocessing Steps}
\begin{itemize}
\item \textbf{Missing value handling:} None required; the dataset was
verified to contain zero missing cells.
\item \textbf{Standardization:} \texttt{StandardScaler} was fit on the
training split only and applied to both Logistic Regression and every SVM
kernel, since both the logistic loss's regularization term and the SVM
margin are scale-sensitive.
\item \textbf{Train--test split:} \texttt{train\_test\_split} with
\texttt{test\_size=0.2}, \texttt{stratify=y}, \texttt{random\_state=42}.
\end{itemize}

\section{Exploratory Data Analysis}
The reusable function \texttt{generate\_eda\_summary()} from
\texttt{ml\_lab\_utils.py} (Experiment 1, unmodified) was used to produce
the consolidated 12-panel EDA figure below.

\begin{figure}[H]
\centering
\includegraphics[width=\linewidth]{figures/eda_spambase.png}
\caption{Consolidated 12-panel EDA summary for the Spambase dataset.}
\label{fig:eda}
\end{figure}

\textbf{Inference:} As established in Experiment 2, the dataset has no
missing values, a moderate 60.6\%/39.4\% class imbalance, and its two
highest-variance features (\texttt{capital\_run\_length\_total} and
\texttt{capital\_run\_length\_longest}) are extremely right-skewed. This
strong non-normality and the resulting need for feature standardization
directly motivates using \texttt{StandardScaler} before both Logistic
Regression and every SVM kernel in this experiment, since both model
families assume or benefit from comparably-scaled inputs.

\section{Implementation Details}

\textbf{Logistic Regression} models the probability that an email is spam
using the sigmoid function,
$P(y{=}1\mid\mathbf{x}) = \frac{1}{1+e^{-(\mathbf{w}^\top\mathbf{x}+b)}}$,
and is trained by minimizing log-loss (cross-entropy). A threshold of 0.5 is
applied to the predicted probability to obtain the class label.
Regularization controls overfitting: \textbf{L1 (Lasso)} shrinks some
coefficients exactly to zero (feature selection), while \textbf{L2 (Ridge)}
shrinks all coefficients toward zero without eliminating any, generally
improving generalization. The inverse regularization strength $C$ trades
off model complexity against regularization: small $C$ means strong
regularization (simpler model), large $C$ means weak regularization
(more complex model, closer to unregularized logistic regression).

\textbf{Support Vector Machine (SVM)} finds the hyperplane that maximizes
the margin between the two classes; only the support vectors (points
closest to the margin) determine the decision boundary. Four kernels were
evaluated: \textbf{Linear} (best for linearly separable data),
\textbf{Polynomial} (captures polynomial feature interactions),
\textbf{RBF} (handles complex non-linear boundaries via a Gaussian
similarity measure), and \textbf{Sigmoid} (a neural-network-activation-like
kernel). SVM's $C$ controls the margin-width/misclassification trade-off
(small $C$: wider margin, higher bias; large $C$: narrower margin, lower
bias), and $\gamma$ controls how far the influence of a single training
point reaches (only relevant for RBF, polynomial, and sigmoid kernels).

\begin{lstlisting}[language=Python, caption={Baseline Logistic Regression via the reusable classification function, from \texttt{ml\_lab\_utils.py} (Experiment 1).}]
baseline_models = {"Logistic Regression (Baseline)": LogisticRegression(max_iter=2000)}
baseline_results_df, baseline_fitted = train_evaluate_classification(
    baseline_models, X_train_scaled, X_test_scaled, y_train, y_test, scale=False)
\end{lstlisting}

\begin{lstlisting}[language=Python, caption={SVM kernel comparison and hyperparameter search space, from \texttt{experiment4\_spambase\_logreg\_svm.py}.}]
kernels = ["linear", "poly", "rbf", "sigmoid"]
for kernel in kernels:
    svm = SVC(kernel=kernel, random_state=42)
    svm.fit(X_train_scaled, y_train)

svm_param_grid = [
    {"kernel": ["linear"], "C": [0.1, 1, 10, 100]},
    {"kernel": ["rbf"], "C": [0.1, 1, 10, 100], "gamma": ["scale", "auto"]},
    {"kernel": ["sigmoid"], "C": [0.1, 1, 10, 100], "gamma": ["scale", "auto"]},
    {"kernel": ["poly"], "C": [0.1, 1, 10], "gamma": ["scale"], "degree": [2, 3]},
]
svm_grid = GridSearchCV(SVC(random_state=42), svm_param_grid,
                          cv=StratifiedKFold(5, shuffle=True, random_state=42),
                          scoring="accuracy")
\end{lstlisting}

\emph{Note on search space:} the polynomial-kernel grid was restricted to
$C\in\{0.1,1,10\}$, $\gamma=\texttt{scale}$, degree $\in\{2,3\}$ (rather than
the full manual-specified grid) to keep total tuning time tractable on the
single-core lab environment used for this run; this does not affect the
selected overall-best kernel, since RBF was already the GridSearchCV winner
by a clear margin.

\section{Experimental Setup}
\begin{tabular}{ll}
\toprule
Python Version & 3.12 \\
Scikit-Learn Version & 1.8.0 \\
IDE & Jupyter Notebook / VS Code \\
Operating System & Ubuntu 24.04 LTS \\
Hardware & Standard lab workstation (CPU only) \\
\bottomrule
\end{tabular}

\vspace{0.3cm}
\textbf{Environment activation command (as mandated by the lab manual):}
\begin{lstlisting}[language=bash]
source /opt/anaconda3/bin/activate anaconda-navigation
\end{lstlisting}

\section{Hyperparameter Tuning Results}

\begin{tabular}{lllc}
\toprule
Model & Search Method & Best Parameters & Best CV Accuracy \\
\midrule
Logistic Regression & Grid / Random & $C=100$, L1, liblinear      & 0.9264 \\
SVM                  & Grid / Random & RBF, $C=10$, $\gamma$=scale & \textbf{0.9356} \\
\bottomrule
\end{tabular}

\vspace{0.2cm}
\noindent GridSearchCV and RandomizedSearchCV agreed on identical (or
functionally identical) best parameters and CV accuracy for both models:
Logistic Regression converged to $C=100$ with L1 penalty and the
\texttt{liblinear} solver (CV accuracy 0.9264 vs. 0.9261 for grid vs.\
random respectively -- effectively tied), and SVM converged to the exact
same configuration for both search methods (RBF, $C=10$, $\gamma$=scale,
CV accuracy 0.9356 for both). RandomizedSearchCV was noticeably faster for
both models (Logistic Regression: 31.6\,s vs.\ 58.6\,s; SVM: 73.7\,s vs.\
83.3\,s) while reaching the same quality solution.

\section{Logistic Regression Performance}
\begin{tabular}{lc}
\toprule
Metric & Value \\
\midrule
Accuracy         & 0.9251 \\
Precision        & 0.9250 \\
Recall           & 0.9251 \\
F1 Score         & 0.9248 \\
Training Time (s)& 0.847 \\
\bottomrule
\end{tabular}

\begin{figure}[H]
\centering
\includegraphics[width=0.7\linewidth]{figures/cm_logreg.png}
\caption{Confusion matrix for the tuned Logistic Regression model on the held-out test set.}
\label{fig:cm_lr}
\end{figure}
\textbf{Inference:} The tuned model achieves balanced precision and recall,
indicating it does not systematically favor one class over the other. Most
residual errors are borderline emails whose word/character frequency
profile sits between typical spam and ham patterns.

\section{SVM Kernel-wise Performance}
\begin{tabular}{lccc}
\toprule
Kernel & Accuracy & F1 Score & Training Time (s) \\
\midrule
Linear  & 0.9294 & 0.9093 & 0.387 \\
Polynomial & 0.7796 & 0.6220 & 0.348 \\
RBF     & \textbf{0.9273} & \textbf{0.9055} & 0.186 \\
Sigmoid & 0.8849 & 0.8528 & 0.243 \\
\bottomrule
\end{tabular}

\begin{figure}[H]
\centering
\includegraphics[width=0.75\linewidth]{figures/svm_kernel_accuracy.png}
\caption{SVM kernel comparison (default hyperparameters) by test accuracy.}
\label{fig:kernel_acc}
\end{figure}
\textbf{Inference:} With default hyperparameters, the Linear and RBF kernels
perform almost identically and far better than Sigmoid, and dramatically
better than the (untuned) Polynomial kernel, whose default degree-3
polynomial expansion appears to distort the decision boundary badly on this
57-dimensional standardized feature space. This gap motivates the
hyperparameter search in Section 6, which found that a properly tuned RBF
kernel ($C=10$) achieves the best cross-validated accuracy of all
configurations tested.

\begin{figure}[H]
\centering
\includegraphics[width=\linewidth]{figures/svm_decision_boundaries_pca.png}
\caption{Decision boundaries of all four SVM kernels, visualized on a 2D PCA projection of the training data.}
\label{fig:decision_boundaries}
\end{figure}
\textbf{Inference:} Projected onto the first two principal components, the
Linear and Polynomial kernels produce very similar, nearly straight decision
boundaries. The RBF kernel wraps a tighter, curved region around the dense
ham cluster, reflecting its ability to model local, non-linear structure.
The Sigmoid kernel produces a visibly distorted, non-convex boundary that
misclassifies large regions of the spam cluster (bottom-left), which is
consistent with its lower accuracy in Table 3 above -- the sigmoid kernel is
generally less numerically stable and more sensitive to hyperparameter
choice than RBF on standardized data.

\begin{figure}[H]
\centering
\includegraphics[width=0.7\linewidth]{figures/cm_svm.png}
\caption{Confusion matrix for the tuned SVM (RBF, $C=10$) model on the held-out test set.}
\label{fig:cm_svm}
\end{figure}
\textbf{Inference:} The tuned SVM's confusion matrix is similarly balanced
to the Logistic Regression model's, with a comparable overall error rate;
the two models make similar types of errors, which is consistent with their
near-identical ROC-AUC scores (Section 9).

\section{K-Fold Cross-Validation Results (K = 5)}

\begin{tabular}{lcc}
\toprule
Fold & Logistic Regression & SVM \\
\midrule
1 & 0.9121 & 0.9305 \\
2 & 0.9359 & 0.9293 \\
3 & 0.9272 & \textbf{0.9413} \\
4 & 0.9283 & 0.9326 \\
5 & 0.9261 & 0.9370 \\
\midrule
Average & 0.9259 & \textbf{0.9341} \\
\bottomrule
\end{tabular}

\begin{figure}[H]
\centering
\includegraphics[width=0.75\linewidth]{figures/cv_accuracy.png}
\caption{5-fold cross-validation accuracy: Logistic Regression vs. SVM.}
\label{fig:cv}
\end{figure}
\textbf{Inference:} SVM's cross-validation accuracy is both higher on
average (0.9341 vs.\ 0.9259) and comparably stable across folds relative to
Logistic Regression, giving confidence that SVM's advantage in cross-
validated accuracy is not an artifact of a single lucky fold.

\section{Comparative Analysis}

\begin{figure}[H]
\centering
\includegraphics[width=\linewidth]{figures/roc_curves.png}
\caption{ROC curves for tuned Logistic Regression and tuned SVM.}
\label{fig:roc}
\end{figure}
\textbf{Inference:} The two models are nearly indistinguishable on ROC-AUC
(Logistic Regression 0.9678 vs.\ SVM 0.9702), with SVM holding a very slight
edge, particularly at low false-positive rates -- a desirable property for a
spam filter that must avoid misclassifying legitimate (ham) email.

\begin{figure}[H]
\centering
\includegraphics[width=\linewidth]{figures/model_comparison_bar.png}
\caption{Logistic Regression vs. SVM: Accuracy, Precision, Recall, and F1 comparison.}
\label{fig:bar}
\end{figure}
\textbf{Inference:} On the single 80/20 test split, tuned Logistic Regression
edges out tuned SVM on accuracy (0.9251 vs.\ 0.9207), precision, recall, and
F1 by a small margin; however, the 5-fold cross-validated results
(Section~7) show the opposite ranking, with SVM ahead on average accuracy.
This apparent discrepancy is a useful illustration of why relying on a
single train/test split can be misleading -- the cross-validated numbers are
the more trustworthy estimate of true generalization performance for both
models.

\begin{figure}[H]
\centering
\includegraphics[width=\linewidth]{figures/regularization_effect.png}
\caption{Left: Logistic Regression test accuracy vs. $C$, by penalty type. Right: number of non-zero coefficients vs. $C$ (L1 sparsity effect).}
\label{fig:reg_effect}
\end{figure}
\textbf{Inference:} Accuracy for both L1 and L2 peaks around $C=1$ and
degrades at the extremes -- very small $C$ (heavy regularization,
$C=0.01$) noticeably hurts L1 more than L2 (0.870 vs.\ 0.908 accuracy),
since L1 is aggressively zeroing out features that still carry useful
signal at that strength. The sparsity plot confirms this directly: at
$C=0.01$, L1 retains only 26 of 57 features, rising to the full 57 by
$C=1$ and beyond, while L2 never zeroes out any feature at any tested $C$.
This is the clearest empirical demonstration in this experiment of L1's
feature-selection behavior versus L2's pure shrinkage behavior.

\begin{tabular}{lcc}
\toprule
Criterion & Logistic Regression & SVM \\
\midrule
Accuracy (test split) & \textbf{0.9251} & 0.9207 \\
Accuracy (5-fold CV)  & 0.9259 & \textbf{0.9341} \\
Model Complexity      & Low & High \\
Training Time         & Low (0.85\,s) & Low--Moderate (0.99\,s) \\
Interpretability      & High & Low \\
\bottomrule
\end{tabular}

\section{Observations}
\begin{itemize}
\item \textbf{Best-performing classifier:} By cross-validated accuracy
(the more reliable estimate), the tuned \textbf{SVM with an RBF kernel}
is the best-performing classifier (0.9341 vs.\ 0.9259), though the margin
over tuned Logistic Regression is modest and the two models are
statistically close on ROC-AUC (0.970 vs.\ 0.968).
\item \textbf{Impact of regularization:} L1 regularization in Logistic
Regression performs automatic feature selection, discarding up to 31 of 57
features at strong regularization ($C=0.01$) with a real accuracy cost;
L2 never discards a feature and is more robust at strong regularization.
Both converge to near-identical, unregularization-level accuracy once
$C\geq 1$, confirming that this dataset does not require heavy
regularization to avoid overfitting -- 4{,}601 samples is ample for a
57-feature linear model.
\item \textbf{Kernel behavior in SVM:} Linear and RBF kernels are the clear
winners; the polynomial kernel performs poorly with its default degree of
3 (likely overfitting to noise or producing numerically unstable feature
expansions in 57 dimensions), and the sigmoid kernel's boundary
(Figure~\ref{fig:decision_boundaries}) is visibly less well-suited to this
feature space's geometry than RBF's tighter, more localized boundary.
\item \textbf{Bias--variance trade-off:} Logistic Regression, being a
simpler linear model, exhibits comparatively higher bias, expressed here as
its plateau in accuracy across a wide range of $C$ values (Table in
Section~9) rather than a sharp drop-off, indicating it is not overfitting
even at low regularization. SVM's RBF kernel, being more flexible, is more
sensitive to hyperparameter choice ($C$ and $\gamma$) -- the untuned RBF SVM
(default $C=1$) already performed well (0.9273 accuracy), but tuning $C$ up
to 10 further reduced bias and improved cross-validated accuracy to 0.9356,
illustrating the classic trade-off of higher model flexibility (lower bias,
if properly tuned) against greater sensitivity to hyperparameter
misspecification.
\end{itemize}

\section{Conclusion}
This experiment implemented and tuned Logistic Regression and SVM
classifiers for spam detection on the Spambase dataset, evaluating four SVM
kernels and comparing L1 vs.\ L2 regularization in Logistic Regression.
GridSearchCV and RandomizedSearchCV converged on essentially identical
optimal hyperparameters for both models (Logistic Regression: $C=100$, L1;
SVM: RBF kernel, $C=10$), with RandomizedSearchCV consistently faster.
Tuned Logistic Regression achieved marginally higher single-split test
accuracy (0.9251), while tuned SVM (RBF) achieved higher and more reliable
5-fold cross-validated accuracy (0.9341) and a slightly higher ROC-AUC
(0.9702). Given SVM's cross-validated advantage, it is the recommended
classifier when raw predictive accuracy is the priority; Logistic
Regression remains an attractive alternative when model interpretability
(explicit, signed feature coefficients) or the lowest possible training
cost is a priority, since it trains and predicts nearly as fast as SVM
while remaining directly interpretable. Both models substantially
outperform the untuned Polynomial and Sigmoid SVM kernels, reinforcing that
kernel choice is at least as important as regularization strength when
building an SVM-based spam classifier.

\section{References}
\begin{enumerate}
\item Scikit-learn Developers, ``Logistic Regression,'' scikit-learn
documentation. [Online]. Available:
\url{https://scikit-learn.org/stable/modules/linear_model.html#logistic-regression}
\item Scikit-learn Developers, ``Support Vector Machines,'' scikit-learn
documentation. [Online]. Available:
\url{https://scikit-learn.org/stable/modules/svm.html}
\item Scikit-learn Developers, ``Tuning the hyper-parameters of an
estimator,'' scikit-learn documentation. [Online]. Available:
\url{https://scikit-learn.org/stable/modules/grid_search.html}
\item M. Hopkins, E. Reeber, G. Forman, and J. Suermondt, ``Spambase Data
Set,'' UCI Machine Learning Repository, 1999. [Online]. Available:
\url{https://archive.ics.uci.edu/dataset/94/spambase}
\end{enumerate}

\end{document}
